\documentclass[11pt,a4paper]{article}

\usepackage[margin=2.5cm]{geometry}
\usepackage{amsmath}
\usepackage{amssymb}
\usepackage{graphicx}
\usepackage{booktabs}
\usepackage{float}
\usepackage{caption}
\usepackage[colorlinks=true, linkcolor=blue, citecolor=blue, urlcolor=blue]{hyperref}
\usepackage{microtype}
\usepackage{parskip}

\title{\textbf{Deep Q-Network for ATARI}\\[0.5em]
       \large Reinforcement Learning Assignment 1 --- BIU 2026}
\author{Yonatan Kapach --- 318515756}
\date{May 2026}

\begin{document}

\maketitle
\tableofcontents
\newpage

% =========================================================================
\section{Game Description}
% =========================================================================

The chosen game is \textbf{Crazy Climber} (Atari 2600), accessed via the
Arcade Learning Environment (ALE) through the Gymnasium interface
(\texttt{ALE/CrazyClimber-v5}).
In Crazy Climber the player controls a man climbing a skyscraper; points
are awarded for ascending each floor and for reaching the top of each
building.
The game emits strictly non-negative rewards, with common per-step bonuses
ranging from 100 to 400 points.
The key challenge for an RL agent is the long horizon required to reach
high floors and accumulate large rewards.

Published benchmark scores for Crazy Climber:
\begin{itemize}
  \item Random agent: 10{,}781
  \item Best linear features agent: 23{,}411
  \item DQN paper: 114{,}103
\end{itemize}

% =========================================================================
\section{Preprocessing}
% =========================================================================

The preprocessing pipeline exactly reproduces Section~4.1 of Mnih et al.\
(2013), implemented using standard Gymnasium wrappers with no external RL
library logic. We used (\texttt{ALE/CrazyClimber-v5}) instead of the required (\texttt{ALE/CrazyClimber-NoFrameskip-v4})
simply because the NoFrameskip version was not available in the Gymnasium registry at the time of implementation. The
parameters of the environment were set to match the required NoFrameskip-v4 ones. 

\begin{enumerate}
  \item \textbf{No-op reset.}
        At the start of each episode a random number of no-operation
        actions (uniformly sampled from $[1, 30]$) are taken to introduce
        stochasticity in the initial state (\texttt{noop\_max=30}).
  \item \textbf{Frame skipping.}
        Each selected action is repeated for 4 consecutive frames; the
        observation returned is the pixel-wise maximum over the last 2
        frames to remove flickering artefacts (\texttt{frame\_skip=4}).
  \item \textbf{Grayscale and resize.}
        Each frame is converted to grayscale and resized to
        $84 \times 84$ pixels using \texttt{AtariPreprocessing}.
  \item \textbf{Frame stacking.}
        The 4 most recent processed frames are stacked along the channel
        dimension, yielding a state tensor of shape $(4, 84, 84)$
        (\texttt{FrameStackObservation}).
  \item \textbf{Sticky actions disabled.}
        \texttt{repeat\_action\_probability=0.0} disables sticky actions,
        matching the NoFrameskip-v4 deterministic-action setting required by
        the assignment.
  \item \textbf{Reward transformation.}
        Applied during training only; raw game scores are used during
        evaluation so that reported rewards are comparable to published
        benchmarks (see Section~\ref{sec:experiments}).
\end{enumerate}

% =========================================================================
\section{Network Architecture}
% =========================================================================

The Q-network follows the convolutional architecture from the DQN paper.
Input observations are stored as \texttt{uint8} in the replay buffer and
normalised to $[0,1]$ (divided by 255) before each forward pass.

\begin{table}[H]
\centering
\caption{DQN network architecture}
\begin{tabular}{llll}
\toprule
Layer   & Type           & Parameters                          & Output shape \\
\midrule
Input   & ---            & ---                                 & $4\times84\times84$ \\
Conv1   & Conv2d + ReLU  & 32 filters, $8\times8$, stride~4    & $32\times20\times20$ \\
Conv2   & Conv2d + ReLU  & 64 filters, $4\times4$, stride~2    & $64\times9\times9$ \\
Conv3   & Conv2d + ReLU  & 64 filters, $3\times3$, stride~1    & $64\times7\times7$ \\
Flatten & ---            & ---                                 & $3{,}136$ \\
FC1     & Linear + ReLU  & 512 units                           & $512$ \\
Output  & Linear         & $n_\text{actions}$ units            & $9$ \\
\bottomrule
\end{tabular}
\end{table}

Crazy Climber exposes 9 discrete actions.
The loss is the Huber loss (smooth-$\ell_1$) applied to temporal-difference
errors computed against a periodically frozen \emph{target network}:
\begin{equation}
  y_i = \tilde{r}_i + \gamma\,\max_{a'} Q_{\theta^{-}}(s_i', a'),
  \qquad
  \mathcal{L}(\theta) =
    \frac{1}{B}\sum_{i=1}^{B}\,\ell_\delta\!\bigl(y_i - Q_\theta(s_i,a_i)\bigr),
\end{equation}
where $\theta^{-}$ denotes the target-network parameters, $\gamma=0.99$,
and $\ell_\delta$ is the Huber loss with $\delta=1$.

% =========================================================================
\section{Hyperparameters}
% =========================================================================

\begin{table}[H]
\centering
\caption{Hyperparameters used for the three final submitted runs}
\begin{tabular}{ll}
\toprule
Hyperparameter                        & Value \\
\midrule
Optimizer                             & Adam \\
Learning rate                         & $1\times10^{-4}$ \\
Discount factor $\gamma$              & $0.99$ \\
Batch size                            & $32$ \\
Replay buffer capacity                & $150{,}000$ transitions \\
Minimum replay size (warm-up)         & $20{,}000$ transitions \\
Network update frequency              & Every $4$ environment steps \\
Target network sync frequency         & Every $10{,}000$ environment steps \\
$\varepsilon$-start (training)        & $1.0$ \\
$\varepsilon$-end (training)          & $0.1$ \\
$\varepsilon$ decay schedule          & Linear over $1{,}000{,}000$ steps \\
$\varepsilon$ during evaluation       & $0.05$ \\
Frame skip                            & $4$ \\
Frame stack ($n$)                     & $4$ \\
Reward mode                           & Clip $\rightarrow \{-1,\,0,\,+1\}$ \\
Total training steps                  & $2{,}000{,}000$ (Run~1); $1{,}900{,}000$ (Runs~2--3) \\
Evaluation frequency                  & Every $10{,}000$ training steps \\
Evaluation episodes per checkpoint    & $30$ \\
\bottomrule
\end{tabular}
\end{table}

\paragraph{Replay buffer capacity.}
The DQN paper uses a buffer of $1{,}000{,}000$ transitions.
To remain within the host-RAM limits of a Google Colab GPU instance the
capacity was reduced to $150{,}000$.
The agent still learns effectively at this size; transitions age out faster
but the diversity requirement is met.

% =========================================================================
\section{Experiments: What Worked and What Did Not}
\label{sec:experiments}
% =========================================================================

\subsection{What Worked}

\paragraph{Adam optimizer.}
Switching from RMSProp (the paper's default) to Adam was the single most
impactful change.
Adam's adaptive first- and second-moment estimates produce faster initial
learning and more stable convergence on Crazy Climber.
Adam produced measurable reward improvements within the first 50{,}000 steps.

\paragraph{Reward clipping to $\{-1,\,0,\,+1\}$.}
Clipping rewards normalises the magnitude of TD targets and prevents
Q-value divergence, a critical stabilisation technique from the original
paper.
For Crazy Climber, which emits only non-negative rewards, reward clipping
simplifies to:
\begin{equation}
  \tilde{r} =
  \begin{cases}
    1 & \text{if } r > 0, \\
    0 & \text{if } r = 0.
  \end{cases}
\end{equation}
This provides a uniform, strong gradient signal for every positive
transition, enabling the agent to focus on maximising the visit frequency
of rewarding states rather than on reward magnitude differentiation.

\paragraph{30 evaluation episodes per checkpoint.}
Using 30 episodes per checkpoint reduces the variance of the mean-reward
estimate and produces a more stable selection signal (but still noisy) 
compared to using only a few single episodes.

\subsection{What Did Not Work}

\paragraph{RMSProp optimizer.}
The paper specifies RMSProp with $\alpha=0.95$, $\varepsilon=0.01$, and
$\text{lr}=2.5\times10^{-4}$.
On Crazy Climber this configuration showed no learning within the first
several hundred thousand steps: mean evaluation rewards remained at or
below the random-agent baseline (${\approx}10{,}781$).
The run was terminated early and Adam was adopted in all subsequent
experiments.

\paragraph{Reward normalization.}
As an alternative to hard clipping, a soft \emph{reward normalization}
scheme was tested:
\begin{equation}
  \tilde{r} = \min\!\left(\frac{r}{r_{\max}},\;1\right)
            = \operatorname{clip}\!\left(\frac{r}{400},\;0,\;1\right),
  \label{eq:normalize}
\end{equation}
where $r_{\max} = 400$ is the maximum dense per-step reward in Crazy
Climber (the climbing bonus on Building~4).

The intent was to preserve ordinal reward structure: under
\eqref{eq:normalize},
$r=100 \Rightarrow \tilde{r}=0.25$,
$r=200 \Rightarrow \tilde{r}=0.50$, and
$r=400 \Rightarrow \tilde{r}=1.0$.
In contrast, hard clipping collapses all positive rewards to $\tilde{r}=1$,
discarding magnitude information but providing a uniform gradient signal.

Normalization proved \textbf{less successful} for the following reasons:

\begin{itemize}
  \item \textbf{Weaker early signal.}
        The Bellman target is $y = \tilde{r} + \gamma\max_{a'}Q(s',a')$.
        Under normalization, the most common small rewards ($r\approx100$)
        contribute only $\tilde{r}=0.25$ to the target, compared to
        $\tilde{r}=1$ under clipping --- a $4\times$ reduction in the
        immediate reward component.
        This significantly slows early learning when the Q-values are
        still near zero and the bootstrapped term is uninformative.

  \item \textbf{Non-uniform gradient magnitudes.}
        Different reward magnitudes produce different TD-error scales,
        requiring the critic to simultaneously handle a fine-grained reward
        scale and the bootstrapping signal.
        This introduces additional instability compared to the simpler
        binary signal from clipping.

  \item \textbf{Empirical outcome.}
        The normalization run peaked at a mean evaluation reward of
        approximately $66{,}700$ (with 10 evaluation episodes per
        checkpoint), compared to peaks of ${\sim}92{,}000$--$96{,}000$
        under reward clipping with 30 evaluation episodes.
\end{itemize}

Hard clipping was retained for all three final reported runs.

% =========================================================================
\section{Results}
% =========================================================================

\subsection{Evaluation Protocol}

\begin{itemize}
  \item Evaluation frequency: every $10{,}000$ environment steps.
  \item Episodes per evaluation checkpoint: \textbf{30}.
  \item Evaluation policy: $\varepsilon$-greedy with $\varepsilon=0.05$.
  \item Rewards recorded: raw (unclipped) game scores.
\end{itemize}

\paragraph{Step selection methodology.}
To avoid overfitting the reported score to any individual run, a single
evaluation step was selected \emph{using only Run~1} and then applied
unchangeably to all three runs:
\begin{enumerate}
  \item Run~1 was trained for $2{,}000{,}000$ steps (199 evaluation
        checkpoints).
  \item The checkpoint with the highest mean evaluation reward in Run~1
        was identified as step $1{,}900{,}000$. (There was actually a slightly higher reward 
        in a much earlier step, but I interpreted that as an outlier and chose the more stable later peak)
  \item Runs~2 and~3 were trained for exactly $1{,}900{,}000$ steps
        (the same budget) using different random seeds.
  \item The final score for each run is the mean evaluation reward at
        step $1{,}900{,}000$, and the aggregate result is the average
        of those three values.
\end{enumerate}
This procedure ensures that the reported step is not cherry-picked across
runs, keeping the evaluation protocol unbiased.

\subsection{Learning Curves}

Figures~\ref{fig:lc_r1}--\ref{fig:lc_r3} show the mean evaluation reward
vs.\ training steps for each of the three independent final runs.
Each point corresponds to one evaluation checkpoint (30 episodes,
$\varepsilon=0.05$).

\begin{figure}[H]
  \centering
  \includegraphics[width=\textwidth]{results/config1_adam_fast_final_v4_r1/learning_curve.png}
  \caption{Learning curve --- Run~1 (full $2{,}000{,}000$ steps, used for
           step selection). Mean evaluation reward at step $1{,}900{,}000$:
           $92{,}160$.}
  \label{fig:lc_r1}
\end{figure}

\begin{figure}[H]
  \centering
  \includegraphics[width=\textwidth]{results/config1_adam_fast_final_v4_r2/learning_curve_1.9M.png}
  \caption{Learning curve --- Run~2 ($1{,}900{,}000$ steps).
           Mean evaluation reward at step $1{,}900{,}000$: $68{,}497$.}
  \label{fig:lc_r2}
\end{figure}

\begin{figure}[H]
  \centering
  \includegraphics[width=\textwidth]{results/config1_adam_fast_final_v4_r3/learning_curve_1.9M.png}
  \caption{Learning curve --- Run~3 ($1{,}900{,}000$ steps).
           Mean evaluation reward at step $1{,}900{,}000$: $83{,}243$.}
  \label{fig:lc_r3}
\end{figure}

\subsection{Final Scores}

The reported score for each run is the mean evaluation reward at the
pre-selected step $1{,}900{,}000$ (30 episodes, $\varepsilon=0.05$).

\begin{table}[H]
\centering
\caption{Final results --- three independent runs at step $1{,}900{,}000$}
\begin{tabular}{lcc}
\toprule
Run   & Seed & Mean eval reward at step $1{,}900{,}000$ \\
\midrule
Run 1 & 42   & $92{,}160$ \\
Run 2 & 100  & $68{,}497$ \\
Run 3 & 467  & $83{,}243$ \\
\midrule
\textbf{Average} & --- & $\mathbf{81{,}300}$ \\
\bottomrule
\end{tabular}
\end{table}

For reference: random agent $= 10{,}781$; best linear agent $= 23{,}411$;
DQN paper $= 114{,}103$.
The achieved average of $\mathbf{81{,}300}$ far exceeds the linear agent
target and reaches ${\approx}71\%$ of the DQN paper score in
$1{,}900{,}000$ training steps.
The learning curve seem to continue trending upwards at the end of training,
suggesting that further training could yield even higher scores. I tend to believe that
with better computational resources and more training steps, the agent could reach or even surpass the DQN paper score. 

\end{document}
